\begin{ouverture}[denombrement]{11}{Graphes et matrices d’adjacence}{Graphes et matrices}{Degrés · Connexité · Chaînes · Puissances}
\objectif{Modéliser une situation par un graphe.}
\objectif{Calculer ordres, degrés et nombres d’arêtes.}
\objectif{Passer du graphe à sa matrice d’adjacence.}
\objectif{Compter des parcours de longueur fixée par une puissance.}
\prerequis{Calcul matriciel et dénombrement élémentaire.}
\end{ouverture}
\essentiel

\essSection{1}{Des sommets reliés}
Un graphe possède des sommets et des arêtes. Deux sommets reliés par une arête sont adjacents.
Dans ce chapitre, les premiers graphes sont simples et non orientés : pas de boucle, ni d'arêtes multiples.
L'ordre est le nombre de sommets ; le degré d'un sommet est le nombre d'arêtes qui lui sont incidentes.
La somme des degrés vaut deux fois le nombre d'arêtes, car chaque arête possède deux extrémités.
Le graphe complet $K_n$ relie toutes les paires de sommets : chacun a degré $n-1$ et il y a $n(n-1)/2$ arêtes.
\figurelivre{e11-graphe.pdf}{Graphe d'ordre $4$, avec les arêtes $AB,AC,BC,CD$. Les degrés sont $2,2,3,1$ ; leur somme $8$ compte deux fois les quatre arêtes.}
\essSection{2}{Chaînes et connexité}
Une chaîne est une succession de sommets dont deux consécutifs sont adjacents. Sa longueur est le nombre d'arêtes parcourues.
Un graphe est connexe si toute paire de sommets peut être reliée par une chaîne.
Dans un graphe orienté, les arêtes deviennent des arcs et un chemin doit respecter leur sens. Un poids peut décrire une distance, un coût ou une probabilité, à préciser selon la situation.
Pour les comptages matriciels, on autorise la répétition des sommets et des arêtes : on compte des parcours, et non uniquement des chemins simples sans répétition.
\essSection{3}{La matrice d'adjacence}
On fixe un ordre des sommets. La matrice $A$ possède un coefficient $a_{ij}=1$ si l'on peut aller directement du sommet $i$ au sommet $j$, et $0$ sinon.
Pour un graphe simple non orienté, elle est symétrique, sa diagonale est nulle et la somme d'une ligne est le degré du sommet.
\begin{exemple}
Pour le graphe de la figure, dans l'ordre $A,B,C,D$,
$A=\begin{pmatrix}0&1&1&0\\1&0&1&0\\1&1&0&1\\0&0&1&0\end{pmatrix}$.
Une permutation de l'ordre des sommets change la présentation de la matrice ; il faut conserver les étiquettes pour interpréter les coefficients.
\end{exemple}
\essSection{4}{Compter avec les puissances}
\begin{theoreme}
Le coefficient $(i,j)$ de $A^n$ compte les parcours de longueur $n$ du sommet $i$ au sommet $j$.
Pour $n=0$, $A^0=I$ correspond au parcours vide, possible seulement si le départ et l'arrivée coïncident.
\end{theoreme}
\begin{demonstration}
Le résultat est la définition pour $n=1$. Supposer le comptage connu à la longueur $n$.
Tout parcours de longueur $n+1$ se décompose en un parcours de longueur $n$ allant de $i$ à un sommet intermédiaire $k$, puis un arc ou une arête de $k$ à $j$.
Pour chaque $k$, le nombre vaut $(A^n)_{ik}a_{kj}$ ; les cas sont disjoints. La somme sur $k$ est précisément $(A^{n+1})_{ij}$.
\end{demonstration}
\figurelivre{e11-parcours.pdf}{Le parcours de longueur $2$ de $A$ à $D$ passe par $C$. Il est compté par $(A^2)_{AD}=1$. Les parcours fermés de longueur $2$ de $A$ à lui-même sont $A,B,A$ et $A,C,A$.}
\essSection{5}{Un prolongement : parcourir chaque arête}
Une chaîne eulérienne utilise chaque arête exactement une fois. Un cycle eulérien revient au départ.
Dans un graphe connexe non orienté avec au moins une arête, il existe une chaîne eulérienne si et seulement si le nombre de sommets de degré impair est $0$ ou $2$ (critère admis ici). Si ce nombre est $0$, un cycle existe.
La nécessité s'explique : dans les sommets intermédiaires, chaque entrée s'associe à une sortie. Seuls le départ et l'arrivée peuvent avoir un degré impair.
Dans le graphe du cours, $D,C,A,B,C$ utilise chaque arête une fois, de $D$ à $C$.

\exercices
\markright{Prendre confiance}
\begingroup\renewcommand{\otBesoinPlace}[1]{}\begin{minipage}{\linewidth}
\rubrique{Prendre confiance}
\begin{exercice}[1]{Vocabulaire}
Un graphe possède cinq sommets et sept arêtes.
\begin{questions}
\item Quel est son ordre ?
\item Quelle est la somme de ses degrés ?
\item Peut-on déjà dire qu’il est complet ?
\end{questions}
\end{exercice}
\begin{exercice}[1]{Graphe du cours}

\begin{questions}
\item Donner les voisins de $C$.
\item Donner le degré de $D$.
\item Le graphe est-il connexe ?
\end{questions}
\end{exercice}
\begin{exercice}[1]{Matrice d’un triangle}

\begin{questions}
\item Écrire la matrice d’adjacence de $K_3$.
\item Donner la somme de chaque ligne.
\item Donner le nombre total d’arêtes.
\end{questions}
\end{exercice}
\begin{exercice}[1]{Une longueur}
On parcourt $A,C,B,A,C,D$ dans le graphe du cours.
\begin{questions}
\item Quelle est la longueur ?
\item Le sommet $A$ est-il répété ?
\item Ce parcours est-il eulérien ?
\end{questions}
\end{exercice}
\end{minipage}\endgroup
\clearpage
\markright{Appliquer et relier}
\begingroup\renewcommand{\otBesoinPlace}[1]{}\begin{minipage}{\linewidth}
\rubrique{Appliquer et relier}
\begin{exercice}[2]{Calculer $A^2$}
$A$ est la matrice du graphe du cours.
\begin{questions}
\item Calculer $(A^2)_{AD}$.
\item Calculer $(A^2)_{AA}$.
\item Interpréter ces deux résultats.
\end{questions}
\end{exercice}
\begin{exercice}[2]{Un graphe orienté}
Dans l’ordre $A,B,C$, $M=\begin{pmatrix}0&1&0\\0&0&1\\1&0&0\end{pmatrix}$.
\begin{questions}
\item Décrire les arcs.
\item Calculer $M^3$.
\item Donner le nombre de chemins de longueur $4$ de $A$ à $B$.
\end{questions}
\end{exercice}
\begin{exercice}[2]{Complet}

\begin{questions}
\item Calculer le nombre d’arêtes de $K_6$.
\item Donner les degrés.
\item Peut-on avoir un cycle eulérien dans $K_6$ ?
\end{questions}
\end{exercice}
\begin{exercice}[2]{Une chaîne eulérienne}
On utilise le graphe du cours.
\begin{questions}
\item Repérer les sommets impairs.
\item Donner une chaîne eulérienne.
\item Peut-elle être fermée ?
\end{questions}
\end{exercice}
\end{minipage}\endgroup
\clearpage
\markright{Chercher et argumenter}
\begingroup\renewcommand{\otBesoinPlace}[1]{}\begin{minipage}{\linewidth}
\rubrique{Chercher et argumenter}
\begin{exercice}[3]{Interpréter une puissance}
Pour un graphe orienté, $(A^5)_{23}=17$.
\begin{questions}
\item Que signifie ce nombre ?
\item Garantit-il $17$ parcours sans répétition ?
\item Que compterait $(A^0)_{23}$ ?
\end{questions}
\end{exercice}
\begin{exercice}[3]{Connexité et information}
Une matrice non orientée est composée de deux blocs diagonaux, sans coefficient non nul entre les blocs. Chaque bloc contient au moins un sommet.
\begin{questions}
\item Que révèle l’absence de coefficients croisés ?
\item Le graphe est-il connexe ?
\item Les puissances peuvent-elles créer un parcours entre les blocs ?
\end{questions}
\end{exercice}
\begin{exercice}[3]{Démontrer le comptage}

\begin{questions}
\item Expliquer le rôle du dernier sommet intermédiaire $k$.
\item Pourquoi multiplie-t-on $(A^n)_{ik}$ par $a_{kj}$ ?
\item Pourquoi additionne-t-on sur $k$ ?
\end{questions}
\end{exercice}
\begin{exercice}[3]{Poids et comptage}
Une matrice contient des distances sur ses arcs plutôt que des $0$ et des $1$.
\begin{questions}
\item Les coefficients de ses puissances comptent-ils encore les parcours ?
\item Que calcule le produit ordinaire dans ce cas ?
\item Quelle précaution prendre avant de choisir une matrice ?
\end{questions}
\end{exercice}
\end{minipage}\endgroup
